\documentclass[11pt,a4paper]{article}

\usepackage[margin=2.4cm]{geometry}
\usepackage[T1]{fontenc}
\usepackage[utf8]{inputenc}
\usepackage{lmodern}
\usepackage{microtype}
\usepackage{amsmath,amssymb}
\usepackage{graphicx}
\usepackage{subcaption}
\usepackage{booktabs}
\usepackage{tabularx}
\usepackage{enumitem}
\usepackage[hidelinks]{hyperref}

% figure folders, relative to this file
\newcommand{\figone}[1]{../code/task1/results/#1}
\newcommand{\figtwo}[1]{../code/task2/results/#1}
\newcommand{\figthree}[1]{../code/task3/results/#1}

% qualitative examples: plain indented block with a bold heading
\newenvironment{example}[1]{\begin{quote}\begin{minipage}{\linewidth}\small\noindent\textbf{#1}\par\smallskip}{\end{minipage}\end{quote}}
\newcommand{\Q}[1]{\noindent\textit{Question:} #1\par}
\newcommand{\RA}[1]{\noindent\textit{Right:} #1\par}
\newcommand{\HA}[1]{\noindent\textit{Hallucinated:} #1\par}
\newcommand{\Samp}[1]{\noindent\textit{Qwen samples:} #1\par}
% right-aligned score: stays on the answer's last line if it fits, else moves to the next line
\newcommand{\score}[1]{{\unskip\nobreak\hfil\penalty50\hskip1em\hbox{}\nobreak\hfil\mbox{#1}\parfillskip=0pt\par}}

\setlist{itemsep=2pt, topsep=4pt}
\setlength{\parskip}{3pt}
\captionsetup{font=small, labelfont=bf}

\title{\textbf{Do LLMs Know When They're Making Things Up?}\\[4pt]
\large Token probabilities, SelfCheckGPT and a stacked detector on HaluEval QA}
\author{Susanta Khamrui (26CS91R12)\\[2pt] \small Project 7, Group 21}
\date{November 2026}

\begin{document}
\maketitle

\begin{abstract}
\noindent
This report tests three hallucination detectors on 200 questions from HaluEval QA (400 answers,
half hallucinated), using Qwen2.5-3B-Instruct. Mean token log-probability and mean entropy
score below random (AUC-PR 0.41 and 0.42). SelfCheckGPT, run with the authors' code, gives 0.65
for the NLI variant and 0.32 for BERTScore. A logistic regression over all four features reaches
0.951 on a held-out 20\%, but answer length alone reaches 0.965. Hallucinated answers in HaluEval
are about four times longer than right ones, and most of each detector's result follows from
that. With length removed, the combined features still score 0.64 AUC-PR.
\end{abstract}

%=====================================================================
\section{Introduction}

The question behind this project is whether a language model's own signals can tell when its
answer is made up. Task~1 uses token probabilities. Task~2 reproduces
SelfCheckGPT~\cite{manakul2023selfcheckgpt}, which checks whether the model agrees with itself
across several samples. Task~3 combines both in one trained classifier.

The assigned dataset is the QA part of HaluEval~\cite{li2023halueval}. Each question (from
HotpotQA) comes with one right answer and one hallucinated answer written by ChatGPT. The two
answers look very different: right answers are short spans like \emph{``yes''} or
\emph{``Lanoxin''}, while hallucinated answers are usually full sentences that repeat the
question. This difference turns out to explain most of the results.

%=====================================================================
\section{Data and setup}\label{sec:data}

200 of the 10{,}000 HaluEval QA pairs were sampled with \texttt{random.Random(42)}. Each pair
gives two items: the right answer (label~0) and the hallucinated answer (label~1). The classes
are balanced, so random guessing scores AUC-PR $=0.5$. A sample of 200 was used because Task~2
needs 20 generations per question on one GPU.

All tasks use \texttt{Qwen/Qwen2.5-3B-Instruct}~\cite{qwen2025} in float16 on an RTX~5060~Ti
(16\,GB). The prompt is closed-book: a system message (\emph{``Answer the question concisely and
factually.''}) and the question. HaluEval's \texttt{knowledge} paragraph is not given to the
model, and nothing is fine-tuned.

\begin{table}[h]
\centering\small
\begin{tabular}{lcccc}
\toprule
 & Median tokens & Mean tokens & Range & Example \\
\midrule
Right answers (label 0)        & 3.0  & 3.7  & 1--13 & \emph{no} \\
Hallucinated answers (label 1) & 14.5 & 16.1 & 2--46 & \emph{Yes, both documentaries\ldots} \\
\bottomrule
\end{tabular}
\caption{Answer length in tokens. The hallucinated answer is longer in 194 of 200 pairs, and
ranking by length alone gives AUC-PR 0.972.}
\label{tab:length}
\end{table}

Table~\ref{tab:length} shows the length gap. Because of it, answer length is reported as a
reference next to every detector below.

%=====================================================================
\section{Task 1: Probability baseline}\label{sec:task1}

\subsection{Method}

For each item, the question prompt $x$ and the answer $a=(a_1,\dots,a_T)$ go through the model
in one forward pass. At each answer position, the log-probability of the actual token and the
entropy of the next-token distribution are recorded:
\begin{equation}
\bar{\ell}(a) = \frac{1}{T}\sum_{t=1}^{T} \log p_\theta(a_t \mid x, a_{<t}),
\qquad
\bar{H}(a) = \frac{1}{T}\sum_{t=1}^{T} H\!\left(p_\theta(\cdot \mid x, a_{<t})\right).
\end{equation}
The hallucination scores are $-\bar{\ell}$ and $\bar{H}$, as in the specification. The minimum
token log-probability and maximum token entropy were also kept for comparison.

\subsection{Results}

\begin{table}[h]
\centering\small
\begin{tabular}{lcc}
\toprule
Score (higher = more suspicious) & AUC-PR & AUROC \\
\midrule
$-$mean log-prob (spec)       & 0.406 & 0.381 \\
mean entropy (spec)           & 0.417 & 0.419 \\
$-$min token log-prob         & 0.660 & 0.697 \\
max token entropy             & 0.610 & 0.643 \\
\midrule
Random                         & 0.500 & 0.500 \\
Answer length only             & 0.972 & 0.971 \\
\bottomrule
\end{tabular}
\caption{Task 1 on all 400 items.}
\label{tab:task1}
\end{table}

\begin{figure}[t]
\centering
\begin{subfigure}{0.49\textwidth}\includegraphics[width=\linewidth]{\figone{pr_curves.png}}\caption{Precision--recall}\end{subfigure}\hfill
\begin{subfigure}{0.49\textwidth}\includegraphics[width=\linewidth]{\figone{roc_curves.png}}\caption{ROC}\end{subfigure}
\caption{Task 1 curves. The two specified scores (solid) fall below the random line; the
worst-token scores (dotted) are above it.}
\label{fig:task1curves}
\end{figure}

Both specified scores are below random (Table~\ref{tab:task1}, Figure~\ref{fig:task1curves}).
In 63\% of pairs the model gives the hallucinated answer a higher mean log-probability than the
right one. Median mean log-probability is $-3.62$ for right answers and $-2.49$ for
hallucinated ones.

The reason is easy to see in a single pair:

\begin{example}{Pair 1290: the model knows the answer, but log-prob prefers the hallucination}
\Q{Do the documentaries Traceroute and War/Dance have the same director?}
\RA{no \score{mean log-prob $-16.30$ (1 token)}}
\HA{Yes, both documentaries were directed by Johannes Grenzfurthner. \score{mean log-prob $-3.48$ (14 tokens)}}
\Samp{``No, Traceroute and War/Dance do not have the same director\ldots'' (all 20 samples start with ``No'')}
\end{example}

The model clearly knows the answer, but it would write \emph{``No, \ldots''} as a full sentence.
A bare lowercase \emph{``no''} is unlikely as a first token, and with only one token nothing
averages it out. The hallucinated answer starts by repeating the question, so most of its tokens
are easy to predict. Mean log-probability here measures style more than truth. The worst-token
scores do better (0.66) because an invented name is often a single surprising token, though they
are also helped by length, since longer answers have more chances to contain a rare token.

%=====================================================================
\section{Task 2: Reproducing SelfCheckGPT}\label{sec:task2}

\subsection{Method}

HaluEval has no pre-generated samples, so $N=20$ answers per question were generated with the
same model and prompt as Task~1 (temperature 1.0, top-$p=1$, no top-$k$, at most 48 new tokens,
seed $42+\text{pair\_id}$). Both answers of a pair are scored against the same samples.

Scoring uses the authors' package \texttt{selfcheckgpt} 0.1.7. Each answer is passed as one
sentence, with the 20 samples as \texttt{sampled\_passages}:
\begin{align}
S_{\text{BERT}}(r) &= 1 - \frac{1}{N}\sum_{n=1}^{N} \max_{k}\; \text{BERTScore-F1}\big(r,\, s^{n}_{k}\big), \\
S_{\text{NLI}}(r)  &= \frac{1}{N}\sum_{n=1}^{N} P\big(\text{contradiction} \mid r,\, S^{n}\big),
\end{align}
where $s^{n}_{k}$ are the sentences of sample $S^{n}$. BERTScore~\cite{zhang2020bertscore} uses
\texttt{roberta-large}~\cite{liu2019roberta} with baseline rescaling. The NLI model is
DeBERTa-v3-large~\cite{he2023debertav3} trained on MNLI~\cite{williams2018mnli}.

\paragraph{Deviations.}
\begin{itemize}
  \item \texttt{microsoft/deberta-v3-large-mnli}, named in the specification, does not exist on
        the Hugging Face Hub. \texttt{potsawee/deberta-v3-large-mnli} is used instead; it is the
        checkpoint the official code loads.
  \item \texttt{selfcheckgpt} 0.1.7 calls \texttt{batch\_encode\_plus}, which transformers~5
        removed. A short shim (\texttt{code/task2/compat.py}) restores it. With the shim, the
        package reproduces the example scores in its README to six decimal places.
  \item The BERTScore variant ignores sample sentences of three tokens or fewer. For two
        questions all 20 samples were single words (e.g.\ ``Canadian''), so those 4 items have no
        BERTScore ($n=396$).
\end{itemize}

\subsection{Results}

\begin{table}[h]
\centering\small
\begin{tabular}{lccccc}
\toprule
 & AUC-PR & AUROC & \multicolumn{2}{c}{Median score} & Pairs ranked \\
 & & & right & halluc. & correctly \\
\midrule
SelfCheck-BERTScore ($n=396$) & 0.317 & 0.075 & 0.959 & 0.571 & 6\% \\
SelfCheck-NLI ($n=400$)       & 0.653 & 0.614 & 0.920 & 0.966 & 67\% \\
\midrule
Random                         & 0.500 & 0.500 & --- & --- & 50\% \\
\bottomrule
\end{tabular}
\caption{Task 2 on all items. ``Pairs ranked correctly'': share of questions where the
hallucinated answer gets the higher score.}
\label{tab:task2}
\end{table}

\begin{figure}[t]
\centering
\begin{subfigure}{0.49\textwidth}\includegraphics[width=\linewidth]{\figtwo{pr_curves.png}}\caption{Precision--recall}\end{subfigure}\hfill
\begin{subfigure}{0.49\textwidth}\includegraphics[width=\linewidth]{\figtwo{roc_curves.png}}\caption{ROC}\end{subfigure}
\caption{Task 2 curves. NLI is above random; BERTScore is far below it, i.e.\ consistent but
in the wrong direction.}
\label{fig:task2curves}
\end{figure}

SelfCheck-NLI beats random (Table~\ref{tab:task2}, Figure~\ref{fig:task2curves}), with good
precision at the top of the ranking: at 20\% recall, 83\% of flagged answers are hallucinated.
Still, even right answers have a median contradiction score of 0.92.

SelfCheck-BERTScore is inverted. It scores the right answer as more suspicious in 94\% of pairs;
with the sign flipped it would get 0.918. The cause is again length:

\begin{example}{Pair 117: why BERTScore runs backwards}
\Q{Jake David Shapiro was best know for being the screenwriter of a film directed by who?}
\RA{Roger Christian \score{$S_{\text{BERT}} = 1.10$}}
\HA{Jake David Shapiro was best known for being the screenwriter of a film directed by Steven Spielberg. \score{$S_{\text{BERT}} = 0.34$}}
\Samp{``Jake David Shapiro was best known for being the screenwriter of the film `Superman II' which was directed by Richard Donner.''}
\end{example}

The samples repeat the question, so a hallucinated answer that also repeats it overlaps with
them heavily. A two-word answer has low recall against any full sentence, whatever it says.

\subsection{When self-consistency works}

SelfCheckGPT assumes the model's samples are mostly right. For each question, the number of
samples containing the right answer was counted. Only 41\% of questions have even one such
sample, and only 20\% (39 questions) have at least 10 of 20.

\begin{figure}[t]
\centering
\includegraphics[width=0.7\textwidth]{\figtwo{auc_by_knowledge.png}}
\caption{AUC-PR split by whether at least 10 of the 20 samples contain the right answer. With
a cut-off of 5 the known-question NLI score is 0.89; with 15 it is 0.93.}
\label{fig:knowledge}
\end{figure}

On the 39 questions the model knows, SelfCheck-NLI reaches \textbf{0.91}; on the other 161 it
falls to 0.59 (Figure~\ref{fig:knowledge}). Leaving out the 12 yes/no questions among the 39
still gives 0.90. Two failure patterns appear on unknown questions:

\begin{example}{Pair 1169: the model is consistently wrong}
\Q{Which film director directed the internet video serial ``God, Inc.'', Roman Karmen or Francis Stokes?}
\RA{Francis Stokes \score{$S_{\text{NLI}} = 0.999$}}
\HA{Roman Karmen directed ``God, Inc.'' with Francis Stokes. \score{$S_{\text{NLI}} = 0.965$}}
\Samp{``Roman Karmen directed the internet video serial `God, Inc.''' (all 20 samples name Roman Karmen)}
\end{example}

\begin{example}{Pair 9763: the model does not know}
\Q{In which city did this world games take place when Squash was added as part of the world games?}
\RA{Lahti, Finland \score{$S_{\text{NLI}} = 0.999$}}
\HA{The World Games was never held in Lahti. \score{$S_{\text{NLI}} = 0.990$}}
\Samp{``Toronto''; ``Riga, Latvia''; ``Doha, Qatar''; ``\ldots Kuala Lumpur, Malaysia in 2017.''}
\end{example}

In pair 1169 the samples agree with each other but not with the truth, so the right answer is
flagged. SelfCheckGPT measures agreement with the model, not with facts. In pair 9763 the
samples scatter, so both answers are contradicted and nothing separates them.

\subsection{Comparison with the paper}

The paper reports results on WikiBio, where 73\% of sentences are inaccurate, so random already
scores 72.96. Table~\ref{tab:paper} therefore also shows the gain over random.

\begin{table}[h]
\centering\small
\begin{tabular}{lcccc}
\toprule
 & \multicolumn{2}{c}{Paper (WikiBio, NonFact)} & \multicolumn{2}{c}{This work (HaluEval QA)} \\
 \cmidrule(lr){2-3}\cmidrule(lr){4-5}
Method & AUC-PR & vs.\ random & AUC-PR & vs.\ random \\
\midrule
Random                  & 72.96 & ---     & 50.0 & --- \\
Avg($-\log p$)          & 83.21 & $+10.3$ & 40.6 & $-9.4$ \\
SelfCheck-BERTScore     & 81.96 & $+9.0$  & 31.7 & $-18.3$ \\
SelfCheck-NLI           & 92.50 & $+19.5$ & 65.3 & $+15.3$ \\
SelfCheck-NLI, known questions & --- & --- & 91.0 & $+41.0$ \\
\bottomrule
\end{tabular}
\caption{Reproduction vs.\ the numbers in the official SelfCheckGPT repository (AUC-PR $\times 100$).}
\label{tab:paper}
\end{table}

NLI's gain over random (+15) is close to the paper's (+19.5), and NLI beats BERTScore as in the
paper. The other rows do not reproduce. Likely causes:
\begin{enumerate}
  \item \textbf{Label source.} In WikiBio, the scored text was written by the model being
        tested. In HaluEval, ChatGPT wrote the hallucinations, so Qwen has no reason to be unsure
        about them, and they differ in style from the right answers.
  \item \textbf{Model knowledge.} GPT-3 (175B) on biographies vs.\ a 3B model on multi-hop
        questions it mostly cannot answer. On questions it does know, NLI matches the paper.
  \item \textbf{Scoring unit.} The paper compares sentences with similar sentences; here
        one-word answers are compared with full-sentence samples, which hurts BERTScore most.
  \item \textbf{Scale.} 400 items give wide intervals (about $\pm0.1$ for NLI on the test split).
\end{enumerate}

%=====================================================================
\section{Task 3: Stacked ensemble}\label{sec:task3}

\subsection{Method}

The 200 \emph{questions} were split 80/20 with \texttt{train\_test\_split(random\_state=42)}
(320 training items, 80 test items). Splitting by question keeps both answers of a question in
the same part. The four features ($\bar{\ell}$, $\bar{H}$, $S_{\text{BERT}}$, $S_{\text{NLI}}$)
are standardised and fed to a logistic regression (L2, $C=1$). The four missing BERTScore values
are filled with the training median.

All detectors are evaluated on the same 80 test items. Task~1 and Task~2 are shown both as raw
scores and as a logistic regression on their own two features; the second is the fairer
comparison. 95\% intervals come from 1000 bootstrap resamples of the 40 test questions.

\subsection{Results}

\begin{table}[h]
\centering\small
\begin{tabular}{lccc}
\toprule
Detector (held-out 20\%, $n=80$) & AUC-PR & 95\% CI & AUROC \\
\midrule
Task 1: $-$mean log-prob             & 0.429 & 0.40--0.51 & 0.422 \\
Task 1: mean entropy                 & 0.459 & 0.42--0.54 & 0.491 \\
Task 1: LR(log-prob, entropy)        & 0.602 & 0.53--0.72 & 0.676 \\
\midrule
Task 2: SelfCheck-BERTScore          & 0.323 & 0.32--0.34 & 0.075 \\
Task 2: SelfCheck-NLI                & 0.651 & 0.56--0.75 & 0.602 \\
Task 2: LR(BERTScore, NLI)           & 0.957 & 0.91--0.99 & 0.959 \\
\midrule
\textbf{Task 3: LR on all four features} & \textbf{0.951} & \textbf{0.90--0.98} & \textbf{0.951} \\
Task 3+: also min log-prob, max entropy  & 0.973 & 0.95--1.00 & 0.972 \\
\midrule
Answer length only                   & 0.965 & 0.94--0.99 & 0.969 \\
Random                               & 0.500 & ---        & 0.500 \\
\bottomrule
\end{tabular}
\caption{All detectors on the same held-out 80 items.}
\label{tab:task3}
\end{table}

\begin{figure}[t]
\centering
\begin{subfigure}{0.49\textwidth}\includegraphics[width=\linewidth]{\figthree{pr_curves.png}}\caption{Precision--recall}\end{subfigure}\hfill
\begin{subfigure}{0.49\textwidth}\includegraphics[width=\linewidth]{\figthree{roc_curves.png}}\caption{ROC}\end{subfigure}
\caption{Held-out curves for each task's trained detector, with answer length as a reference.}
\label{fig:task3curves}
\end{figure}

Task~3 reaches 0.951 (Table~\ref{tab:task3}, Figure~\ref{fig:task3curves}), far above the
specified Task~1 and Task~2 scores. It ranks 38 of 40 test pairs correctly and labels 87.5\% of
items correctly at a 0.5 threshold. However, it is level with Task~2's own regression (0.957)
and just below length alone (0.965); the intervals overlap. Over 100 random splits the picture
is the same: LR(Task~1) $0.627 \pm 0.047$, LR(Task~2) $0.948 \pm 0.025$, Task~3
$0.948 \pm 0.025$, length $0.972 \pm 0.014$.

\begin{figure}[t]
\centering
\includegraphics[width=0.65\textwidth]{\figthree{coefficients.png}}
\caption{Standardised weights of the Task 3 regression. Positive pushes towards ``hallucinated''.}
\label{fig:coef}
\end{figure}

The weights explain this (Figure~\ref{fig:coef}). The largest is BERTScore at $-3.08$: the
classifier uses BERTScore in reverse, as Section~\ref{sec:task2} suggests. BERTScore correlates
with log answer length at $r=-0.70$, so most of the ensemble's signal is length.

\subsection{Controlling for length}

Two checks, both fitted on training data only:
\begin{enumerate}
  \item \textbf{Length as a feature:} compare LR on $\log(\text{tokens})$ alone with LR on
        $\log(\text{tokens})$ plus the four features.
  \item \textbf{Length removed:} regress each feature on $\log(\text{tokens})$, keep the
        residuals, and train the regression on those.
\end{enumerate}

\begin{figure}[t]
\centering
\includegraphics[width=0.82\textwidth]{\figthree{length_control.png}}
\caption{Mean AUC-PR over 100 question-level splits (bars: $\pm$1 sd). Grey: features as they
are. Blue: length removed. Orange: length added as a feature.}
\label{fig:lengthcontrol}
\end{figure}

\begin{table}[h]
\centering\small
\begin{tabular}{lcc}
\toprule
Detector (100 splits) & AUC-PR & Pairs correct \\
\midrule
LR(length)                          & $0.972 \pm 0.014$ & 97\% \\
LR(length + Task 1 + Task 2)        & $0.987 \pm 0.008$ & 98\% \\
\midrule
LR(Task 1, length removed)          & $0.616 \pm 0.059$ & 74\% \\
LR(Task 2, length removed)          & $0.667 \pm 0.053$ & 71\% \\
LR(Task 1 + Task 2, length removed) & $0.642 \pm 0.056$ & 75\% \\
$-$mean log-prob alone, length removed & $0.614 \pm 0.063$ & 75\% \\
\bottomrule
\end{tabular}
\caption{Length controls, mean $\pm$ sd over 100 splits.}
\label{tab:lengthcontrol}
\end{table}

The features do carry information beyond length (Figure~\ref{fig:lengthcontrol},
Table~\ref{tab:lengthcontrol}). Added to length, they raise AUC-PR from 0.972 to 0.987 and help
on 95 of 100 splits. With length removed, 0.64 AUC-PR remains and three out of four pairs are
ranked correctly.

Mean log-probability changes the most. On its own it was below random, but with length removed
it ranks 75\% of pairs correctly. Length had been pushing it the wrong way. NLI goes the other
direction: after the correction it is no better than random (0.44 on the seed-42 split), so part
of its raw advantage was length too.

\subsection{Error analysis}

The ensemble mis-orders 2 of the 40 test pairs. In both, the ``hallucinated'' answer is actually
correct. In pair 6572, \emph{``\ldots is The Chipmunks''} is a paraphrase of \emph{``Alvin and
the Chipmunks''}. Pair 3657 is shown below; for this question all 20 Qwen samples say
``Canadian'', so the model itself is wrong.

\begin{example}{Pair 3657: a label problem in HaluEval}
\Q{\ldots songs by Heart, a rock band that first found success in Canada and later in the United States and worldwide, are of which nationality?}
\RA{American \score{ensemble $0.74$}}
\HA{Heart, the famous rock band, is American. \score{ensemble $0.62$}}
\end{example}

There are five false negatives at the 0.5 threshold. Three are specific but wrong (e.g.\ ``Cancun,
Mexico'' instead of Daytona Beach), one is close to a paraphrase, and one makes no checkable
claim:

\begin{example}{Pair 2060: a vague answer nothing contradicts}
\Q{Turbidimetric inhibition immunoassay is used in many commercial immunoassays, including this medication sold under which brand name?}
\RA{Lanoxin \score{ensemble $0.01$, $S_{\text{NLI}} = 0.93$}}
\HA{Turbidimetric inhibition immunoassay is used in many commercial immunoassays, including several medications. \score{ensemble $0.37$, $S_{\text{NLI}} = 0.32$}}
\end{example}

The answer avoids the question instead of inventing a fact, so the samples cannot contradict it.
Self-consistency methods cannot catch answers that say nothing.

%=====================================================================
\section{Discussion}

\paragraph{What worked.} SelfCheck-NLI works when its assumption holds: 0.91 AUC-PR on questions
the model knows, with no training and no access to model internals. Splitting by question and
reporting a length-only baseline made it clear which results are meaningful.

\paragraph{What did not.} The two specified Task~1 scores are below random and SelfCheck-BERTScore
is inverted. The Task~3 ensemble looks strong, but mostly because it learns length through
BERTScore.

\paragraph{Why.} Two reasons. First, the hallucinations were written by a different model in a
different style, so they are not tied to Qwen's uncertainty, which all these methods rely on.
Second, a 3B model answering multi-hop questions closed-book usually does not know the answer,
and self-consistency gives no signal (or the wrong one) in that case.

\paragraph{Limitations.} Only 200 of 10{,}000 questions were used, so intervals are wide. The
``known question'' split uses string matching and misses paraphrases. The length correction is
linear and may not remove all of it. A length-matched test set was not possible: only 28 pairs
have answers within five tokens of each other. A cleaner test would score answers that Qwen
generated itself, labelled against HotpotQA's gold answers.

%=====================================================================
\section{Conclusion}

On HaluEval QA with a 3B open model, token probabilities do not detect hallucinations, and
SelfCheckGPT only works on the questions the model can answer. The stacked regression reaches
0.95 AUC-PR, but almost all of that comes from answer length. With length controlled, the
combined features rank three out of four pairs correctly, which is the more realistic measure of
what they can do.

%=====================================================================
\begin{thebibliography}{9}\small

\bibitem{manakul2023selfcheckgpt}
P.~Manakul, A.~Liusie, and M.~J.~F.~Gales.
SelfCheckGPT: Zero-resource black-box hallucination detection for generative large language models.
In \emph{Proc.\ EMNLP}, 2023. arXiv:2303.08896.
Code: \url{https://github.com/potsawee/selfcheckgpt}.

\bibitem{li2023halueval}
J.~Li, X.~Cheng, W.~X.~Zhao, J.-Y.~Nie, and J.-R.~Wen.
HaluEval: A large-scale hallucination evaluation benchmark for large language models.
In \emph{Proc.\ EMNLP}, 2023. arXiv:2305.11747.

\bibitem{qwen2025}
Qwen Team.
Qwen2.5 technical report.
arXiv:2412.15115, 2024.

\bibitem{zhang2020bertscore}
T.~Zhang, V.~Kishore, F.~Wu, K.~Q.~Weinberger, and Y.~Artzi.
BERTScore: Evaluating text generation with BERT.
In \emph{Proc.\ ICLR}, 2020. arXiv:1904.09675.

\bibitem{liu2019roberta}
Y.~Liu, M.~Ott, N.~Goyal, J.~Du, M.~Joshi, D.~Chen, O.~Levy, M.~Lewis, L.~Zettlemoyer, and V.~Stoyanov.
RoBERTa: A robustly optimized BERT pretraining approach.
arXiv:1907.11692, 2019.

\bibitem{he2023debertav3}
P.~He, J.~Gao, and W.~Chen.
DeBERTaV3: Improving DeBERTa using ELECTRA-style pre-training with gradient-disentangled embedding sharing.
In \emph{Proc.\ ICLR}, 2023. arXiv:2111.09543.

\bibitem{williams2018mnli}
A.~Williams, N.~Nangia, and S.~Bowman.
A broad-coverage challenge corpus for sentence understanding through inference.
In \emph{Proc.\ NAACL}, 2018.

\end{thebibliography}

%=====================================================================
\appendix
\section{Reproducibility}

Code, a pinned \texttt{requirements.txt} and run instructions are in the submitted archive
(\texttt{README.md}). With Docker, everything runs with one command, \texttt{bash run\_docker.sh}
(add \texttt{--recompute} to redo the GPU steps). Without Docker, it is \texttt{bash code/run\_all.sh}.

\begin{table}[h]
\centering\small
\begin{tabularx}{\textwidth}{l>{\raggedright\arraybackslash}X}
\toprule
Item & Setting \\
\midrule
Data & HaluEval \texttt{qa\_data.json} (github.com/RUCAIBox/HaluEval), 200 pairs, \texttt{random.Random(42)} \\
LLM & \texttt{Qwen/Qwen2.5-3B-Instruct}, float16, closed-book chat prompt \\
Resampling & $N=20$, temperature 1.0, top-$p$ 1.0, top-$k$ off, max 48 new tokens, seed $42+\text{pair\_id}$ \\
BERTScore & \texttt{roberta-large} via \texttt{bert-score} 0.3.13, baseline rescaling \\
NLI & \texttt{potsawee/deberta-v3-large-mnli} \\
SelfCheckGPT & \texttt{selfcheckgpt} 0.1.7 (official code, commit \texttt{19b492a}) \\
Split & question-level 80/20, \texttt{train\_test\_split(random\_state=42)} \\
Classifier & \texttt{StandardScaler} + \texttt{LogisticRegression(C=1, random\_state=42)} \\
CIs & 1000 bootstrap resamples of test questions, \texttt{default\_rng(42)} \\
Extra splits & 100 more splits for the length check, \texttt{random\_state} $0,\dots,99$ \\
Hardware & NVIDIA RTX 5060 Ti 16\,GB; Python 3.12, PyTorch 2.14, transformers 5.21 \\
\bottomrule
\end{tabularx}
\end{table}

\end{document}
